% GENERATED FILE. Edit canonical lessons, then run npm run generate:books.
% canonical-source-hash: fnv1a64:848d37b7d959cbc2
% canonical-lessons: ZH-C385-tigao, ZH-C385-bacheng, ZH-C385-jieneng, ZH-C385-zhongshi, ZH-C385-fangshi, ZH-R385-review-380-385

\chapter{Good Reasons}
\label{ch:zh-good-reasons}
\hlchaptermodality{385}

\begin{chapteropening}
\textbf{By the end of this chapter:} \emph{I can talk about raising prices, what is most likely, saving energy, taking things seriously and ways of living.}

Complete the last lesson of chapter 385: I can talk about raising prices, what is most likely, saving energy, taking things seriously and ways of living.
\end{chapteropening}

\section[tígāo]{tígāo — to push up (a level)}

\label{lesson:ZH-C385-tigao}



\begin{quote}
Before the new one: say the Mandarin for natural resources, then the Mandarin for to increase (in number).
\end{quote}

\subsection*{You'll want to know: tígāo}
\emph{tígāo} — \textquotedblleft{}to push up (a level)\textquotedblright{}.

\begin{quote}
\emph{zhè jiā fàndiàn tígāo le jiàqián} — This restaurant has put its prices up.
\end{quote}

Its characters come later; for now it is a word for the ear.

\begin{cousinweb}
\emph{tígāo} is built from \emph{tí} 'to carry by the handle' and \zh{高} 'high', both inherited from Middle Chinese, \emph{dej} and \emph{kaw} (Guangyun, in Baxter's transcription), which glosses the first 'to carry in the hand'.
\end{cousinweb}

\begin{grammarlens}[title={What you've built}]
One more word for giving reasons.
\end{grammarlens}

\subsection*{Your turn}
\begin{itemize}
\raggedright
  \item \emph{Say it:} \emph{tígāo}
  \item \emph{Say it:} \emph{tígāo}, once more
  \item \emph{Say it:} \emph{zhè jiā fàndiàn tígāo le jiàqián} — This restaurant has put its prices up.
  \item \emph{Recall:} say the Mandarin for natural resources, then the Mandarin for to increase (in number), then say \emph{tígāo} again
\end{itemize}

\begin{tcolorbox}[breakable,colback=teal!4,colframe=teal!35!black,title={Before you move on}]
What does tígāo mean? (\textquotedblleft{}To push up (a level)\textquotedblright{}.) Say it once more.
\end{tcolorbox}
\section[\texorpdfstring{\zh{八成} (bāchéng)}{八成 (bāchéng)}]{\zh{八成} (bāchéng) — most likely; eighty per cent}

\label{lesson:ZH-C385-bacheng}



\begin{quote}
Before the new one: say the Mandarin for to increase (in number), then the Mandarin for to push up (a level).
\end{quote}

\subsection*{You'll want to know: \zh{八成}}
\textbf{\zh{八成}} — \emph{bāchéng} — \textquotedblleft{}most likely; eighty per cent\textquotedblright{}.

\begin{quote}
\textbf{\zh{他八成病了所以没来}} \emph{tā bāchéng bìng le suǒyǐ méi lái} — He's most likely ill, which is why he didn't come.
\end{quote}

\begin{cousinweb}
\textbf{\zh{八成}} (\emph{bāchéng}) is built from \zh{八} 'eight' and \zh{成} 'a tenth', both inherited from Middle Chinese, \emph{peat} and \emph{dzyeng} (Guangyun, in Baxter's transcription). The dictionary glosses the second 'to complete'; \zh{成} for one part in ten is later and is not given there.
\end{cousinweb}

\begin{grammarlens}[title={What you've built}]
One more word for giving reasons.
\end{grammarlens}

\subsection*{Your turn}
\begin{itemize}
\raggedright
  \item \emph{Say it:} \emph{bāchéng}
  \item \emph{Say it:} \emph{bāchéng}, once more
  \item \emph{Say it:} \emph{tā bāchéng bìng le suǒyǐ méi lái} — He's most likely ill, which is why he didn't come.
  \item \emph{Recall:} say the Mandarin for to increase (in number), then the Mandarin for to push up (a level), then say \emph{bāchéng} again
\end{itemize}

\begin{tcolorbox}[breakable,colback=teal!4,colframe=teal!35!black,title={Before you move on}]
What does \zh{八成} mean? (\textquotedblleft{}Most likely; eighty per cent\textquotedblright{}.) Say it once more.
\end{tcolorbox}
\section[\texorpdfstring{\zh{节能} (jiénéng)}{节能 (jiénéng)}]{\zh{节能} (jiénéng) — to save energy; energy-saving}

\label{lesson:ZH-C385-jieneng}



\begin{quote}
Before the new one: say the Mandarin for to push up (a level), then the Mandarin for most likely (eight in ten).
\end{quote}

\subsection*{You'll want to know: \zh{节能}}
\textbf{\zh{节能}} — \emph{jiénéng} — \textquotedblleft{}to save energy; energy-saving\textquotedblright{}.

\begin{quote}
\textbf{\zh{我觉得坐火车比开车节能}} \emph{wǒ juéde zuò huǒchē bǐ kāichē jiénéng} — I think taking the train saves more energy than driving.
\end{quote}

\begin{cousinweb}
\textbf{\zh{节能}} (\emph{jiénéng}) is built from \zh{节} 'to restrain' and \zh{能} 'energy', both inherited from Middle Chinese, \emph{tset} and \emph{nong} (Guangyun, in Baxter's transcription), which glosses the first 'the joint of a bamboo' and 'to hold back'. The word is modern.
\end{cousinweb}

\begin{grammarlens}[title={What you've built}]
One more word for giving reasons.
\end{grammarlens}

\subsection*{Your turn}
\begin{itemize}
\raggedright
  \item \emph{Say it:} \emph{jiénéng}
  \item \emph{Say it:} \emph{jiénéng}, once more
  \item \emph{Say it:} \emph{wǒ juéde zuò huǒchē bǐ kāichē jiénéng} — I think taking the train saves more energy than driving.
  \item \emph{Recall:} say the Mandarin for to push up (a level), then the Mandarin for most likely (eight in ten), then say \emph{jiénéng} again
\end{itemize}

\begin{tcolorbox}[breakable,colback=teal!4,colframe=teal!35!black,title={Before you move on}]
What does \zh{节能} mean? (\textquotedblleft{}To save energy; energy-saving\textquotedblright{}.) Say it once more.
\end{tcolorbox}
\section[zhòngshì]{zhòngshì — to attach importance to}

\label{lesson:ZH-C385-zhongshi}



\begin{quote}
Before the new one: say the Mandarin for most likely (eight in ten), then the Mandarin for to save energy.
\end{quote}

\subsection*{You'll want to know: zhòngshì}
\emph{zhòngshì} — \textquotedblleft{}to attach importance to\textquotedblright{}.

\begin{quote}
\emph{zhèngfǔ yīnggāi zhòngshì huánjìng} — The government should take the environment seriously.
\end{quote}

Its characters come later; for now it is a word for the ear.

\begin{cousinweb}
\emph{zhòngshì} is built from \zh{重} 'heavy' and \emph{shì} 'to look at', both inherited from Middle Chinese, \emph{drjowngX} and \emph{dzyijH} (Guangyun, in Baxter's transcription), which glosses the first 'thick' and 'careful' and the second 'to look'.
\end{cousinweb}

\begin{grammarlens}[title={What you've built}]
One more word for giving reasons.
\end{grammarlens}

\subsection*{Your turn}
\begin{itemize}
\raggedright
  \item \emph{Say it:} \emph{zhòngshì}
  \item \emph{Say it:} \emph{zhòngshì}, once more
  \item \emph{Say it:} \emph{zhèngfǔ yīnggāi zhòngshì huánjìng} — The government should take the environment seriously.
  \item \emph{Recall:} say the Mandarin for most likely (eight in ten), then the Mandarin for to save energy, then say \emph{zhòngshì} again
\end{itemize}

\begin{tcolorbox}[breakable,colback=teal!4,colframe=teal!35!black,title={Before you move on}]
What does zhòngshì mean? (\textquotedblleft{}To attach importance to\textquotedblright{}.) Say it once more.
\end{tcolorbox}
\section[fāngshì]{fāngshì — a way of doing things; a mode}

\label{lesson:ZH-C385-fangshi}



\begin{quote}
Before the new one: say the Mandarin for to save energy, then the Mandarin for to attach importance to.
\end{quote}

\subsection*{You'll want to know: fāngshì}
\emph{fāngshì} — \textquotedblleft{}a way of doing things; a mode\textquotedblright{}.

\begin{quote}
\emph{qí zìxíngchē shì yí ge hěn huánbǎo de fāngshì} — Cycling is a very green way to get about.
\end{quote}

Its characters come later; for now it is a word for the ear.

\begin{cousinweb}
\emph{fāngshì} is built from \emph{fāng} 'a method' and \emph{shì} 'a pattern', both inherited from Middle Chinese, \emph{pjang} and \emph{syik} (Guangyun, in Baxter's transcription), which glosses the first 'square' and 'a method' and the second 'a rule'.
\end{cousinweb}

\begin{grammarlens}[title={What you've built}]
One more word for giving reasons.
\end{grammarlens}

\subsection*{Your turn}
\begin{itemize}
\raggedright
  \item \emph{Say it:} \emph{fāngshì}
  \item \emph{Say it:} \emph{fāngshì}, once more
  \item \emph{Say it:} \emph{qí zìxíngchē shì yí ge hěn huánbǎo de fāngshì} — Cycling is a very green way to get about.
  \item \emph{Recall:} say the Mandarin for to save energy, then the Mandarin for to attach importance to, then say \emph{fāngshì} again
\end{itemize}

\begin{tcolorbox}[breakable,colback=teal!4,colframe=teal!35!black,title={Before you move on}]
What does fāngshì mean? (\textquotedblleft{}A way of doing things; a mode\textquotedblright{}.) Say it once more.
\end{tcolorbox}
\section[(dialogue)]{Looking back: from Clean Air to Good Reasons}

\label{lesson:ZH-R385-review-380-385}



\begin{quote}
Say the Mandarin for the environment and the air we breathe.
\end{quote}

\subsection*{Your turn}
\begin{itemize}
\raggedright
  \item \emph{Say it:} the environment, the air we breathe, to pollute, our planet, eco-friendly (good for the environment)
  \item \emph{Say it:} to sort into kinds, the whole family, to recycle, nationwide, a sense of security
  \item \emph{Say it:} to be on duty, a security guard, a security camera, to guarantee, to patrol
  \item \emph{Say it:} the electricity bill, to waste, sources of energy, to cut down on, disposable (used once)
  \item \emph{Say it:} to plant trees, to become, production cost, natural resources, to increase (in number)
  \item \emph{Say it:} to push up (a level), most likely (eight in ten), to save energy, to attach importance to, a way of doing things
  \item \emph{Read it:} \zh{全} \zh{保} \zh{成} — each one aloud, then a word that holds it
\end{itemize}

\begin{tcolorbox}[breakable,colback=teal!4,colframe=teal!35!black,title={Before you move on}]
The environment? (\emph{huánjìng}.) The air we breathe? (\emph{kōngqì}.) To pollute? (\emph{wūrǎn}.) Our planet? (\emph{dìqiú}.) Eco-friendly (good for the environment)? (\emph{huánbǎo}.) To sort into kinds? (\emph{fēnlèi}.) The whole family? (\textbf{\zh{全家}}, \emph{quánjiā}.) To recycle? (\emph{huíshōu}.) Nationwide? (\textbf{\zh{全国}}, \emph{quánguó}.) A sense of security? (\textbf{\zh{安全感}}, \emph{ānquángǎn}.) To be on duty? (\emph{zhíbān}.) A security guard? (\textbf{\zh{保安}}, \emph{bǎo'ān}.) A security camera? (\emph{jiānkòng}.) To guarantee? (\emph{bǎozhèng}.) To patrol? (\emph{xúnluó}.) The electricity bill? (\emph{diànfèi}.) To waste? (\emph{làngfèi}.) Sources of energy? (\emph{néngyuán}.) To cut down on? (\emph{jiǎnshǎo}.) Disposable (used once)? (\emph{yícìxìng}.) To plant trees? (\emph{zhòngshù}.) To become? (\textbf{\zh{成为}}, \emph{chéngwéi}.) Production cost? (\textbf{\zh{成本}}, \emph{chéngběn}.) Natural resources? (\emph{zīyuán}.) To increase (in number)? (\emph{zēngjiā}.) To push up (a level)? (\emph{tígāo}.) Most likely (eight in ten)? (\textbf{\zh{八成}}, \emph{bāchéng}.) To save energy? (\textbf{\zh{节能}}, \emph{jiénéng}.) To attach importance to? (\emph{zhòngshì}.) A way of doing things? (\emph{fāngshì}.)
\end{tcolorbox}
